% !TeX root=../main.tex
% دستور زیر باید در نخستین فصل وجود داشته باشد.
\pagenumbering{arabic}
\thispagestyle{empty}
\chapter{مقدمه}
\label{ch:introduction}

\section{پیش‌گفتار}

توانایی برقراری ارتباط یکی از پایه‌های رشد شناختی، عاطفی و اجتماعی کودک است. کودک پیش از آنکه نخستین واژه‌های خود را بیان کند، با نگاه‌کردن، لبخندزدن، تولید آوا، اشاره و دنبال‌کردن جهت نگاه دیگران در ارتباط شرکت می‌کند. این رفتارهای پیش‌زبانی به‌تدریج با درک واژه‌ها، نامیدن اشیا و ساخت عبارت‌های کوتاه پیوند می‌خورند. سه سال نخست زندگی از فشرده‌ترین دوره‌های اکتساب گفتار و زبان است و وجود محیطی سرشار از گفت‌وگو، صدا، تصویر و تعامل پاسخ‌گو به شکل‌گیری این توانایی‌ها کمک می‌کند \cite{NIDCDMilestones2022}.

رشد زبان صرفاً با شنیدن منفعلانه‌ی واژه‌ها اتفاق نمی‌افتد. کودک هنگامی بهتر می‌آموزد که واژه در یک موقعیت معنادار ارائه شود، فرصت توجه به شیء یا تصویر مرتبط را داشته باشد، بتواند نوبت خود را در تعامل تجربه کند و پاسخ او از سوی بزرگسال دیده و دنبال شود. به همین دلیل، فعالیت‌هایی مانند نام‌بردن اشیای روزمره، کتاب‌خوانی مشترک، بازی هدایت‌شده و تکرار واژه در موقعیت‌های واقعی، بخش مهمی از محیط زبانی کودک را تشکیل می‌دهند. پژوهش‌های مرورشده درباره مداخله‌های اجراشده توسط والد نیز نشان داده‌اند که آموزش راهبردهای حمایت زبانی به والدین با پیامدهای مثبت در ارتباط و زبان کودکان خردسال همراه است \cite{RobertsCurtis2019,RobertsKaiser2011}.

در سال‌های اخیر تلفن همراه و تبلت به بخشی از محیط زندگی خانواده‌ها تبدیل شده‌اند. این ابزارها می‌توانند تصویر، صوت، بازخورد فوری و ثبت روند تمرین را در یک محیط یکپارچه فراهم کنند؛ با این حال، برای کودک زیر سه سال نباید جای تعامل انسانی، بازی حرکتی و تجربه مستقیم با محیط را بگیرند. کیفیت محتوا، سادگی تعامل، همراهی والد و محدودبودن زمان استفاده در این گروه سنی اهمیت ویژه دارد \cite{WHO2019ScreenTime,AAPMedia2016}. پروژه حاضر با درنظرگرفتن همین فرصت‌ها و محدودیت‌ها، بر توسعه برنامه کاربردی «نوواژه» تمرکز دارد.

\section{بیان مسئله}

والدین معمولاً نخستین و پایدارترین همراهان کودک در مسیر یادگیری زبان هستند، اما تبدیل توصیه‌های کلی مانند «با کودک بیشتر صحبت کنید» به تمرین‌هایی منظم و متناسب با علایق هر کودک همیشه آسان نیست. محتوای عمومی موجود ممکن است با زبان خانواده، محیط فرهنگی، اشیای آشنا یا واژه‌هایی که کودک در زندگی روزمره به آن‌ها نیاز دارد هماهنگ نباشد. از سوی دیگر، تکرار بدون تنوع می‌تواند انگیزه کودک را کاهش دهد و والد نیز ابزار روشنی برای مشاهده روند پاسخ‌گویی کودک در طول زمان در اختیار نداشته باشد.

مسئله اصلی این پروژه، طراحی بستری است که ارائه‌ی دیداری و شنیداری واژه، تمرین انتخاب و یادآوری، بازی‌های کوتاه، شخصی‌سازی محتوا توسط والد و ثبت روند یادگیری را کنار هم قرار دهد. در چنین بستری، کودک صدای یک واژه را می‌شنود، تصویر متناظر را پیدا می‌کند و نتیجه انتخاب خود را از طریق بازخورد ساده و قابل‌فهم دریافت می‌کند. والد نیز می‌تواند برای هر فرزند نمایه جداگانه بسازد، واژه‌ها و دسته‌بندی‌های متناسب با محیط زندگی او را اضافه کند و پاسخ‌های ثبت‌شده را مشاهده نماید.

چالش طراحی فقط افزودن تعدادی تصویر و صدا به یک برنامه نیست. مخاطب نهایی هنوز توانایی خواندن ندارد، دامنه توجه او کوتاه است و کنترل حرکتی انگشتانش در حال رشد است. بنابراین رابط کاربری باید از متن و کنترل‌های پیچیده دوری کند، هدف هر صفحه را واضح نگه دارد، گزینه‌های محدود و بزرگ نمایش دهد و از بازخورد سریع بهره ببرد. همچنین لازم است بخش مدیریتی والد از محیط بازی کودک جدا باشد و برای جلوگیری از استفاده طولانی یا خروج ناخواسته از برنامه، کنترل زمان و حالت کودک پیش‌بینی شود.

چالش دیگر، تفسیر درست داده‌های حاصل از تمرین است. پاسخ درست به یک پرسش در برنامه معادل تسلط کامل بر یک واژه یا توانایی کاربرد آن در گفت‌وگوی واقعی نیست. کودک ممکن است تصویر را بر اساس رنگ، جای گزینه یا حدس انتخاب کند و انتقال آموخته از صفحه به محیط واقعی نیز به تعامل والد وابسته است. ازاین‌رو، شاخص‌هایی مانند تعداد تلاش، درصد پاسخ درست و وضعیت «آموخته‌شده» در نوواژه صرفاً نشانه‌های درون‌برنامه‌ای برای تنظیم تمرین هستند و نباید به‌عنوان تشخیص بالینی یا سنجش جامع رشد گفتار تفسیر شوند.

\section{هدف پروژه}

هدف کلی پروژه، همکاری در طراحی و توسعه یک برنامه کاربردی چندسکویی برای کمک به ایجاد فرصت‌های منظم، جذاب و قابل‌شخصی‌سازیِ یادگیری واژه برای کودکان زیر سه سال است. این هدف با تمرکز بر تعامل والد و کودک دنبال می‌شود و برنامه به‌عنوان ابزار کمکی آموزشی در نظر گرفته شده است، نه جایگزین گفتاردرمانگر، ارزیابی شنوایی یا تعامل طبیعی خانواده.

اهداف جزئی پروژه عبارت‌اند از:
\begin{itemize}
    \item ارائه واژه‌های فارسی در دسته‌بندی‌های موضوعی همراه با تصویر و تلفظ صوتی؛
    \item ایجاد تمرین‌های کوتاه مبتنی بر شنیدن واژه و انتخاب تصویر متناظر؛
    \item اولویت‌دادن به واژه‌هایی که بر اساس پاسخ‌های قبلی هنوز تثبیت نشده‌اند؛
    \item افزایش مشارکت کودک با استفاده از بازخورد مثبت و بازی‌های ساده و متنوع؛
    \item فراهم‌کردن امکان ثبت تصویر و صدای آشنا توسط والد برای ساخت محتوای شخصی؛
    \item نگهداری تاریخچه تلاش‌ها و محاسبه شاخص‌های قابل‌فهم برای مرور پیشرفت؛
    \item جداسازی فضای مدیریتی والد از فضای ساده‌شده کودک؛
    \item کنترل مدت استفاده و کاهش احتمال خروج ناخواسته کودک از محیط برنامه؛
    \item اشتراک بخش عمده کد میان اندروید و آی‌اواس برای کاهش هزینه نگهداری و توسعه.
\end{itemize}

\section{راهکار ارائه‌شده: نوواژه}

نوواژه یک سامانه کارخواه‌ـ‌کارساز است که از برنامه همراه، خدمات کارساز و پایگاه داده تشکیل می‌شود. محتوای پایه سامانه مجموعه‌ای از واژه‌های روزمره در گروه‌هایی مانند حیوانات، میوه‌ها، رنگ‌ها، اعداد، وسایل نقلیه، پوشاک، خوراکی‌ها، شکل‌ها و سبزیجات است. هر واژه می‌تواند نام فارسی و انگلیسی، تصویر، فایل صوتی و ترتیب نمایش داشته باشد. در تجربه اصلی یادگیری، تلفظ یک واژه پخش می‌شود و کودک از میان چند تصویر، گزینه متناظر را انتخاب می‌کند. پاسخ به کارساز ارسال و در صورت وجود نمایه فعال برای کودک، نتیجه در سابقه او ذخیره می‌شود.

پرسش‌ها تنها به‌صورت تصادفی تولید نمی‌شوند؛ سامانه واژه‌هایی را که هنوز در وضعیت آموخته‌شده قرار نگرفته‌اند در اولویت قرار می‌دهد. این سازوکار شکل ساده‌ای از تمرین تطبیقی است و مانع می‌شود تمام واژه‌ها بدون توجه به عملکرد پیشین کودک با فراوانی یکسان تکرار شوند. پس از هر پاسخ، بازخورد دیداری و شنیداری فوری نمایش داده می‌شود. افزون بر آزمون شنیداری‌ـ‌تصویری، بازی‌هایی برای حافظه، تطبیق سایه، تشخیص گزینه متفاوت، مرتب‌سازی رنگ و شناخت اجزای چهره در برنامه وجود دارد. این بازی‌ها مستقیماً معادل گفتاردرمانی نیستند، اما می‌توانند تنوع، توجه و آمادگی کودک برای ادامه تعامل را افزایش دهند.

در بخش والد، ثبت‌نام و ورود امن، مدیریت نمایه فرزندان، افزودن دسته‌بندی، مشاهده واژه‌های موجود و ساخت واژه سفارشی فراهم شده است. برای واژه سفارشی، والد نام واژه را وارد می‌کند، تصویر دلخواهی برمی‌گزیند و تلفظ آن را با صدای خود ضبط می‌کند. این ویژگی کمک می‌کند اشخاص، اشیا، خوراکی‌ها یا مکان‌های آشنای خانواده وارد تمرین شوند و فعالیت از یک فهرست عمومی فراتر رود. برنامه همچنین امکان تعیین زمان مجاز بازی و فعال‌سازی حالت کودک را در اختیار والد قرار می‌دهد.

رابط مشترک اندروید و آی‌اواس با کاتلین و کامپوز چندسکویی پیاده‌سازی شده است. منطق ارتباط با کارساز نیز در ماژول مشترک قرار دارد و کارساز مبتنی بر کتور، مسیرهای مدیریت محتوا، احراز هویت، فرزندان، آزمون و پیشرفت را ارائه می‌کند. این انتخاب فنی اجازه می‌دهد رفتار و ظاهر اصلی برنامه در دو سکوی همراه یکپارچه بماند و در عین حال قابلیت‌های وابسته به سکو، مانند ضبط و پخش صوت یا قفل‌کردن برنامه، در لایه‌های ویژه هر سیستم‌عامل پیاده‌سازی شوند.

\section{رویکرد انجام پروژه}

توسعه نوواژه به‌صورت افزایشی و مبتنی بر جداسازی مسئولیت‌ها انجام شده است. ابتدا مدل‌های اصلی دامنه شامل والد، کودک، دسته‌بندی، واژه، پرسش و سابقه پاسخ تعریف شدند. سپس خدمات کارساز و رابط‌های برنامه‌نویسی برای دسترسی به این داده‌ها شکل گرفتند. در سمت برنامه همراه، مخزن‌های داده وظیفه ارتباط با کارساز را بر عهده گرفتند و مدل‌های نمایش، وضعیت هر صفحه را مدیریت کردند. در مرحله بعد، رابط‌های جداگانه والد و کودک، پخش صوت، ضبط صدای والد، انتخاب تصویر، بازی‌ها و بازخوردهای حرکتی و صوتی افزوده شدند.

در تصمیم‌های طراحی، سه اصل دنبال شده است: نخست، انجام هر فعالیت کودک با کمترین تعداد گام؛ دوم، فراهم‌کردن کنترل محتوا و زمان برای والد؛ و سوم، قابل‌اندازه‌گیری‌کردن تعامل بدون تبدیل شاخص‌های نرم‌افزاری به ادعای تشخیصی. نتیجه این رویکرد، نمونه‌ای قابل اجرا و توسعه‌پذیر است که می‌تواند در ادامه با آزمون کاربردپذیری، نظر متخصصان گفتار و زبان و مطالعه میدانی بهبود یابد.

\section{دامنه و محدودیت‌ها}

نوواژه بر یادگیری و مرور واژه و برخی مهارت‌های شناختی پشتیبان تمرکز دارد. نسخه فعلی صدای کودک را برای تشخیص خودکار تلفظ تحلیل نمی‌کند، کیفیت تولید آوا را نمی‌سنجد و برنامه درمانی اختصاصی برای اختلال‌های گفتار یا زبان تولید نمی‌کند. وضعیت آموخته‌شدن واژه از روی تعداد پاسخ‌های درست در آزمون تصویری تعیین می‌شود و بنابراین تنها رفتار کودک در همان فعالیت را بازتاب می‌دهد.

افزون بر این، رشد گفتار تحت تأثیر عوامل گوناگونی مانند شنوایی، رشد شناختی و حرکتی، کیفیت تعامل، چندزبانگی و تفاوت‌های فردی قرار دارد. اگر والد درباره درک زبان، تولید واژه، واکنش کودک به صدا یا پسرفت مهارت‌ها نگرانی داشته باشد، مراجعه به پزشک، شنوایی‌شناس یا آسیب‌شناس گفتار و زبان ضروری است \cite{ASHALateLanguage,NIDCDMilestones2022}. استفاده از برنامه نباید سبب به‌تعویق‌افتادن ارزیابی تخصصی شود.

محدودیت دیگر به سن مخاطب و استفاده از صفحه‌نمایش مربوط است. راهنماهای سلامت عمومی بر فعالیت بدنی، خواب کافی و محدودکردن فعالیت نشسته مبتنی بر صفحه تأکید دارند \cite{WHO2019ScreenTime}. بنابراین الگوی مطلوب استفاده از نوواژه، جلسه‌های کوتاه با حضور والد است؛ والد پس از نمایش تصویر و پخش صدا، همان واژه را در محیط واقعی تکرار می‌کند، به تلاش ارتباطی کودک پاسخ می‌دهد و سپس فعالیت را از صفحه به بازی با اشیای واقعی منتقل می‌سازد.

\section{دستاوردهای پروژه}

مهم‌ترین دستاوردهای این پروژه را می‌توان به‌صورت زیر خلاصه کرد:
\begin{itemize}
    \item تولید یک برنامه همراه فارسی و کودک‌محور با قابلیت اجرا بر اندروید و آی‌اواس؛
    \item یکپارچه‌سازی محتوای دیداری و شنیداری، آزمون، بازی و ثبت پیشرفت در یک سامانه؛
    \item پشتیبانی از محتوای شخصی‌سازی‌شده با تصویر و صدای والد؛
    \item طراحی سازوکار ساده اولویت‌بندی واژه‌های آموخته‌نشده؛
    \item ارائه بخش مدیریتی برای والد و محیط کم‌پیچیدگی برای کودک؛
    \item پیاده‌سازی کارساز مستقل برای احراز هویت، مدیریت محتوا و نگهداری داده‌ها؛
    \item فراهم‌کردن مبنایی برای ارزیابی‌های آتی کاربردپذیری و اثر آموزشی با همکاری متخصصان.
\end{itemize}

\section{ساختار گزارش}

این گزارش در پنج فصل تنظیم شده است. فصل حاضر مسئله، اهداف، دامنه و راهکار کلی نوواژه را معرفی می‌کند. فصل دوم به مفاهیم پایه رشد گفتار و زبان، نقش گفتاردرمانی و والد، یادگیری مبتنی بر بازی، اصول کاربرد رسانه دیجیتال برای کودک خردسال و مبانی مفهومی سامانه می‌پردازد. فصل سوم راهنمای فنی پروژه و جزئیات معماری، اجزا، مدل داده، رابط‌های کارساز و روش اجرای سامانه را ارائه خواهد کرد. فصل چهارم راهنمای استفاده والد و کودک از قابلیت‌های برنامه را در بر می‌گیرد. در فصل پنجم نیز دستاوردها و محدودیت‌ها جمع‌بندی و مسیرهای پیشنهادی برای توسعه و ارزیابی آینده بیان خواهد شد.
